\chapter{与AI对话的艺术：Prompting基础}

第一章讨论了 AI 协作的基本框架：人负责定义目标和边界，AI 负责生成候选答案和执行明确任务。本章进入具体操作层面，回答一个更直接的问题：如何把一个真实需求转化为 AI 能够理解并可靠执行的指令？

这个转化过程并不像“随便问一句”那样简单。即使是同一个问题，不同的提问方式往往导致截然不同的结果。本章介绍 Prompt 的构成要素、核心写作技术和常见陷阱，并说明如何把一次性问答升级为可复用、可管理的协作工作流。

从研究脉络看，Prompt 并不只是向聊天机器人提问的措辞技巧，而是一种通过自然语言组织模型行为的程序接口。“Prompt programming”的概念强调，使用者是在设计任务、上下文和交互过程，而不是寻找一句神奇咒语\cite{reynolds2021prompt}。因此，本章关注的不是孤立的提示技巧，而是如何使提示可解释、可测试、可复用，并能在模型或数据发生变化后被重新验证。

% ----------------------------------------------------------------
\section{从模糊意图到可执行指令}
% ----------------------------------------------------------------

\subsection{信息空白的代价}

与人类专家交流时，很多背景信息可以省略不说：对方知道你们所处的工作环境，了解行业惯例，能从语境中推断出你没有明说的部分。AI 不具备这种共同知识。它接收的信息只有你提供的内容，其余的空白由模型用训练数据中的统计规律填补。这意味着，每一处未被说清楚的细节，都可能变成一个潜在的假设分歧点。

考虑以下两种提问方式：

\begin{quote}
\textbf{提问 A：}帮我分析一下交通流量数据。
\end{quote}

\begin{quote}
\textbf{提问 B：}我有某市 10 个道路检测器 2023 年全年的 5 分钟级数据，字段包含时间戳、检测器编号、车速（km/h）和流量（辆/5 分钟）。其中速度字段中，值为零的记录通常代表检测器故障，而不是车辆实际停止，请不要把它们当作有效的低速记录处理。请完成以下三项分析：（1）统计每个检测器的缺失率和异常零值比例；（2）比较工作日与法定节假日的平均速度分布；（3）用 Python 绘制一张热力图，横轴为一天中的小时（0 到 23），纵轴为检测器编号，颜色表示平均速度，颜色越深代表速度越低。请输出可直接运行的 Python 代码，并在每个分析阶段前添加一行注释说明该步骤的目的。
\end{quote}

提问 A 迫使 AI 猜测：你想分析什么指标？数据包含哪些字段？速度为零应如何处理？输出应该是图表还是文字？结果往往是一段风格通用、对具体数据几乎没有帮助的代码。提问 B 提供了充分信息：数据结构与来源、特殊字段的领域语义、三项具体任务以及输出格式要求。AI 生成直接可用代码的概率显著提高，即使第一次输出仍需调整，调整方向也更加清晰。

\subsection{AI如何填补信息空白}

当 Prompt 缺失关键信息时，模型会自动补充假设。这些假设来自训练数据中最常见的模式，在通用任务中可能恰好适用，但在有特殊约束的专业任务中，往往不适用。

以下是一些典型的自动假设：

\begin{itemize}
  \item 速度为零意味着车辆真实停止，而不是传感器故障；
  \item 训练集和测试集可以按比例随机划分；
  \item“最好的模型”意味着在测试集上使均方误差最小；
  \item 缺失值用列均值填充是合理的默认操作；
  \item 分析中涉及的约束与公开文献中的通用做法一致。
\end{itemize}

这些假设中的每一条，在特定场景下都可能是错误的。数据泄露、不适当的评价指标、忽视领域异常值——这些问题中相当一部分来自信息空白，而不是 AI 能力本身的限制。一个关键推论是：\textbf{Prompt 的质量决定了 AI 所处理问题的质量}。写出更清晰的 Prompt，本质上是在帮助自己更清楚地思考问题本身。

% ----------------------------------------------------------------
\section{Prompt的五个构成要素}
% ----------------------------------------------------------------

一个高质量 Prompt 可以从五个维度来设计：角色、上下文、任务、约束和输出格式。这五个维度并非总是同等重要，也不需要每次都完整填写，但对于专业任务而言，系统地审视这五个问题，往往能够显著减少后续的反复调整。

\subsection{角色（Role）}

角色指令告诉 AI 以什么专业背景和工作立场完成任务。

\begin{quote}
你是一位专注于城市交通安全的工程师，正在为交通管理部门准备月度事故分析报告，读者以行政管理人员为主，无需呈现过多统计细节。
\end{quote}

角色设定主要影响两个方面。其一是输出风格与侧重：要求 AI 扮演“面向管理者的报告作者”和“面向研究人员的数据分析师”，即使面对同一份数据，措辞、详略和关注点也会不同。其二是对未明说的格式和专业惯例的推断：要求 AI 以“公交调度工程师”的角色生成运营方案时，它会更倾向于考虑车辆周转时间、班次衔接和驾驶员排班等约束，而不是只讨论数学模型。

但需要注意：角色设定影响的是输出的风格和优先级，而不是 AI 实际掌握的知识量。声明角色不会让模型拥有它本来不具备的专业知识。如果任务依赖高度专业的细节，仍然需要在上下文中明确提供。

\subsection{上下文（Context）}

上下文是 Prompt 中信息密度最高的部分，也是最容易被低估的部分。它向 AI 解释你所处的具体情况，使 AI 能够在正确的前提下工作。上下文通常包含三个层次：

\begin{description}
  \item[数据上下文] 数据来源、采集方式、字段含义、单位、时间粒度、缺失规律和已知异常。例如：“本数据集来自城市交通管理平台的感应线圈检测器，每 5 分钟汇总一次；speed 字段若超过 200 km/h 属于记录错误，flow 字段值为 -1 代表检测器离线，均应视为缺失值处理。”
  \item[领域上下文] 业务规则、专业惯例与行业约定。例如：“本项目的分析时间范围仅覆盖工作日早高峰（7:00—9:00），不包含法定节假日和夜间时段；'拥堵'定义为连续 15 分钟平均速度低于 20 km/h。”
  \item[项目上下文] 已完成的工作、已确认的决策和当前任务在整体流程中的位置。例如：“上一步已完成数据清洗，输出文件为 clean\_detector.csv。本步骤只需要在这份清洗后的数据上构造时间特征，不需要重复清洗过程。”
\end{description}

提供上下文时，不必追求面面俱到——无关信息会稀释注意力，也可能引入混淆。原则是：只提供 AI 无法自行推断、但对当前任务确实必要的信息。

\subsection{任务（Task）}

任务是 Prompt 的核心，说明 AI 具体需要完成什么。一个清晰的任务应当满足三个条件：

\begin{itemize}
  \item \textbf{具体：}能够明确说出完成标准，而不只是“做个分析”或“建个模型”；
  \item \textbf{有边界：}明确说明任务从哪里开始、到哪里结束，以及不包含哪些不相关的部分；
  \item \textbf{可验证：}输出能够被独立检查，例如代码能运行、表格能对比、结论能用原始来源核实。
\end{itemize}

任务的粒度是关键。“帮我建立一个交通流量预测模型”涵盖的内容太多，无法一次可靠完成；“使用 clean\_detector.csv 中的 flow 字段，计算每个检测器过去 7 天的滚动平均值（以小时为粒度），并将结果作为新列 flow\_7d\_avg 附加到数据集中，输出处理后数据集的前 5 行用于核验”则明确、可执行且可验证。

多个子任务应拆开，按顺序逐一完成，而不是一次性塞入单个 Prompt。每个子任务完成后，应检查输出是否符合预期，再进入下一步。

\subsection{约束（Constraint）}

约束告诉 AI 哪些事情不能做，或必须在特定条件下做。约束是把领域隐含知识转化为明确规则的主要工具。常见约束类型包括：

\begin{description}
  \item[时序约束]“预测任务中，特征只能使用预测时刻之前已经可获得的数据，任何需要未来才能获得的字段均不得使用。”
  \item[数据约束]“不要删除原始数据集中的任何列；所有新增字段写入新列，所有修改后的数据写入新文件，原文件保持不变。”
  \item[技术约束]“代码须兼容 Python 3.10，只使用 pandas、numpy 和 scikit-learn，不引入其他第三方库。”
  \item[领域约束]“事故严重程度的划分使用本地交通管理局的三级分类标准（轻微、一般、严重），不使用通用文献中的分类体系。”
  \item[安全与伦理约束]“本数据集包含个人出行轨迹，分析中不得输出任何能够识别个人身份的信息，包括直接标识和可间接推断的内容。”
\end{description}

约束的价值在于将“默认不说也不应该做”的事情明确化。领域专家往往认为这些约束不言而喻，但对 AI 来说，未被声明的约束等于不存在。数据泄露、不合伦理的隐私处理、不符合运营条件的方案——这些问题中的大部分可以通过主动声明约束来预防，而不是在输出错误之后再纠正。

\subsection{输出格式（Format）}

输出格式规定了 AI 应以什么形式呈现结果。明确指定格式有两个好处：它让输出更容易被检查，也让输出更容易被直接使用。常用格式指定方式包括：

\begin{itemize}
  \item \textbf{代码：}“请输出一段可独立运行的 Python 脚本，文件顶部列出所有 import，每个分析阶段前添加一行注释说明该步骤的目的。”
  \item \textbf{表格：}“请以 Markdown 表格形式输出，包含检测器编号、缺失率、平均速度和异常值比例四列，按缺失率从高到低排序。”
  \item \textbf{分步说明：}“请先列出完整分析步骤，再为每一步生成对应代码，不要直接输出全部代码而不解释步骤划分。”
  \item \textbf{中间推理：}“在给出最终方案之前，请先列出你对问题的理解，以及你考虑过哪些可能的方法，说明为什么选择当前方案，再开始生成代码。”
\end{itemize}

要求输出中间推理过程是一种特别有用的技巧。当 AI 被要求先解释思路再执行时，推理步骤变得可见，使用者可以在执行阶段之前发现方向错误，而不必等到得到完整输出后再回头检查。这在方法选择和数据处理策略等决策节点尤为重要。

% ----------------------------------------------------------------
\section{写出高质量提示的核心技术}
% ----------------------------------------------------------------

掌握五要素之后，还需要了解在不同任务场景下应当采用哪种提示策略。以下介绍四种核心技术，它们并不是互斥的，可以根据任务需要组合使用。

\subsection{零样本与少样本提示}

零样本提示（zero-shot prompting）是指不提供任何示例，只通过描述任务本身让 AI 完成目标。这是最常见的提示形式，适用于任务描述已足够清晰、AI 无需参考具体示例就能正确理解的情况。

少样本提示（few-shot prompting）则是在 Prompt 中提供一至数个示例，明确展示期望的输入输出格式或推理模式。以下情况通常更适合使用少样本提示：

大规模语言模型的研究表明，模型能够仅凭上下文中的少量示例适应新任务\cite{brown2020language}。但示例同时也是隐含规范：它不仅展示格式，还会暗示类别边界、措辞风格和允许的推断方式。因此，少样本示例应覆盖关键边界情况，而不能只挑选最容易、最整齐的样本。否则模型可能很好地模仿示例，却在真实数据的模糊区域稳定失效。

\begin{itemize}
  \item 任务的输出格式非常特殊，难以用文字精确描述；
  \item 任务涉及特定机构的分类规则，与通用标准存在差异；
  \item 输出需要保持高度一致的风格或结构。
\end{itemize}

以交通事件分类为例，如果需要按照某机构的内部标准（而非通用规范）对事件描述进行等级标注，直接描述所有分类规则可能冗长且难以准确把握细节；此时提供几个标注好的示例，往往比用文字解释规则更高效：

\begin{quote}
以下是我们机构使用的事件严重程度分类示例，请按相同规则对新描述进行分类：

\smallskip
\textbf{示例 1}\\
描述：高速公路发生追尾，货车与两辆乘用车碰撞，3 人轻伤，事故车道全部封闭，拥堵延伸至 5 公里以外。\\
分类：A 级（重大拥堵事件）

\smallskip
\textbf{示例 2}\\
描述：城市主干道路缘石旁发现废弃摩托车，交警到场处理，未占用行车道。\\
分类：C 级（轻微道路隐患）

\smallskip
请对以下新描述进行分类：\ldots
\end{quote}

少样本提示中的示例对 AI 的输出具有强烈引导作用。如果示例本身包含错误或偏差，AI 可能学到并延续这些问题。因此，示例的质量与数量同样重要。

\subsection{思维链提示}

思维链提示最初用于通过中间推理步骤改善复杂推理任务表现\cite{wei2022chain}；后续研究还发现，即使不提供示例，简短的分步思考指令也可能提高部分任务的正确率\cite{kojima2022large}。然而，中间推理文本仍然是模型生成的文本，不等于可审计的真实思维过程，更不能替代外部证据。实践中更可靠的做法，是要求 AI 给出\textbf{可检查的中间产物}，例如假设列表、公式、数据查询、测试用例和来源，而不是把一段听起来合理的解释本身视为验证。

思维链提示（chain-of-thought prompting）要求 AI 在给出最终答案之前，先生成中间推理步骤。其核心价值不仅在于“让 AI 想清楚”，更在于让推理过程对人可见，从而能够在结论产生之前发现逻辑错误。

最简单的激活方式是在 Prompt 中加入明确指令，例如“请先列出分析思路，再生成代码”：

\begin{quote}
我想对交通流量序列做短期预测，但不确定应当选择哪类方法。请先列出你认为可以考虑的几类方法，对每种方法说明它依赖的假设、适用的数据结构和主要的局限性；再结合我的数据情况（5 分钟粒度、单个路段数据、有明显日周期规律、有约 8\% 的随机缺失），给出建议方案并说明理由。不要直接跳到最终推荐。
\end{quote}

思维链提示在以下场景中尤其有价值：

\begin{enumerate}
  \item 方法选择类任务：AI 对候选方法的分析步骤可以暴露其对假设的理解，使用者可以发现并纠正不适合当前场景的假设；
  \item 数据处理策略类任务：要求 AI 先解释处理逻辑，再生成代码，可以在执行前发现方向性错误；
  \item 结论可信度需要被理解的任务：当一个结论需要向他人解释时，AI 的推理步骤可以作为起草说明材料的基础。
\end{enumerate}

思维链并不适用于所有任务。对于简单且目标明确的代码生成任务，额外的推理步骤只会增加输出长度而不带来实质价值。判断标准是：跳过推理是否可能导致难以察觉的错误？

\subsection{分步拆解复杂任务}
\label{sec:decomposition}

复杂任务应被分解为多个具有单一目标的子步骤，每个步骤有清楚的输入和可检查的输出。分步执行不只是为了让 AI 更容易完成任务，更是为了让使用者在每一步结束时能够验证结果，避免错误在后续步骤中积累。

以“利用道路检测器数据预测交通流量”为例，可以拆解为以下步骤，每步单独提交 Prompt：

\begin{enumerate}
  \item 读取数据，输出字段列表、数据量、时间范围和基本统计摘要。\textit{验证：字段数量、时间范围与预期一致。}
  \item 检查缺失值、异常速度和异常流量，输出各字段缺失率和异常值比例。\textit{验证：异常比例与领域先验知识相符，没有大量正常数据被误标记。}
  \item 按约定规则处理异常值，输出处理前后的速度和流量对比图。\textit{验证：图形变化符合预期，没有意外地删除大量数据。}
  \item 构造时间特征（小时、星期、节假日标记）和滞后特征，输出特征列表和样例行。\textit{验证：特征不包含未来信息，格式与预期一致。}
  \item 按时间顺序划分训练集与测试集，输出各集合的起止时间和样本量。\textit{验证：测试集完整位于训练集之后，没有时间交叉。}
  \item 训练基准模型，输出训练集和测试集上的误差指标。\textit{验证：基准误差合理，测试集表现不显著优于训练集（否则可能存在数据泄露）。}
\end{enumerate}

注意，每一步的验证标准应在执行前确定，而不是在看到输出之后临时评价。事先确定验证标准，能够有效防止“结果看起来不错就接受”的确认偏误。

\subsection{迭代修正}

第一次输出通常是协作的起点，而不是最终结果。有效的迭代修正建立在清楚分析“第一次输出哪里出了问题”的基础上，而不是含糊地说“不太对，请重试”。

有效的修正 Prompt 应明确说明：具体哪个部分有问题（代码行号、输出格式、错误的处理逻辑）、问题是什么（报错信息、与预期的偏差、被违反的约束），以及需要如何修改（如果已知方向）。

\begin{quote}
你刚才生成的代码在计算 speed\_zero\_ratio 时（第 23 行），以包含缺失值在内的全部记录作为分母。但我要求的是：在非缺失且速度不为零的有效记录中，速度为零的异常比例。请修改这一行，使分母只统计 speed 字段非空且大于零的记录数。
\end{quote}

这比“统计方法有问题，请修改”有效得多。精确的修正使 AI 更容易正确响应，也使对话历史更有条理，便于后续复查。

当一个问题经过两轮以上迭代仍无法解决时，通常意味着需要补充更多上下文，或者任务的定义本身需要重新澄清。继续机械地重复修正往往效率不高；更有价值的做法是退一步，重新审视当前 Prompt 中哪个要素描述得不够清晰。

% ----------------------------------------------------------------
\section{常见提示陷阱与修复方法}
% ----------------------------------------------------------------

理解 Prompt 的构成之后，还需要识别常见的失败模式。以下五类陷阱在实际工作中最为普遍，每类均附有可操作的修复方式。

\subsection{目标模糊，依赖AI猜测}

“帮我做一个交通预测的项目”是这类陷阱的典型代表。这条 Prompt 看似给出了方向，实际上 AI 无从得知预测的是什么指标（流量、速度还是行程时间）、针对哪段道路或区域、预测多远的未来、可用数据是什么、评价标准是什么。AI 只能用训练数据中最常见的假设填补这些空白，结果可能与真实需求完全不符。

修复的关键在于先独立完成问题定义，再动笔写 Prompt。至少需要明确预测对象、时空粒度、可用数据、评价方式和关键约束，将这些内容写入 Prompt 的任务和上下文部分。

\subsection{缺少必要的领域上下文}

“请清洗这份交通数据，处理缺失值和异常值”是另一类常见失败。AI 不知道什么在这份数据中算”异常值”，不知道缺失值应当删除、前向填充还是线性插值，也不知道速度为零是否是有效记录。套用通用数据处理惯例往往会产生错误结果，因为这些惯例在特定领域下并不成立。

修复方式是在上下文中提供字段语义说明、已知数据质量问题以及本领域的处理惯例。例如：

\begin{quote}
本数据来自感应线圈检测器，speed 字段值为零通常代表检测器故障而非车辆真实停止，应将其标记为缺失值，不得保留为零速度记录。缺失值处理策略：连续缺失不超过 3 个时间步时，使用线性插值；超过 3 步的连续缺失段，标记为不可用时段并整体删除。
\end{quote}

\subsection{约束未声明，引发数据泄露或不可用方案}

时序数据建模中有一类危险的隐含约束极易被忽视。以”请用交通数据训练一个预测模型，并报告测试集的 MAE”为例，AI 可能随机划分训练集和测试集而非按时间顺序，将未来信息包含在特征中，或者选择在实际部署时无法获得的数据源作为输入。测试集表现因此可能虚高，而在真实应用中却完全失效。

预防这类问题的方式是明确声明关键约束：

\begin{quote}
训练集和测试集必须按时间顺序划分，使用数据中前 80\% 的时间段作为训练集，后 20\% 作为测试集，禁止随机划分。模型所用特征只能来自预测时刻之前已经可以获得的数据，不能使用任何未来才能获得的字段。
\end{quote}

\subsection{一次要求太多}

“请读取数据、清洗、选择最优特征、训练三种模型、比较结果、输出最终报告”这类 Prompt 把整个分析流程压缩进了一次对话。每个环节都存在独立的判断点，一旦某个环节出错，下游所有步骤的输出都可能失效；而当一次性得到完整输出时，往往很难定位错误出现在哪一步。

解决方式是按第~\ref{sec:decomposition} 节的方法将任务拆解为独立步骤，每步单独发送 Prompt 并验证输出，确认无误后再继续。

\subsection{输出无法验证}

“请为这份交通数据选择最合适的预测模型”看似是一个明确的问题，但 AI 只能给出一个名字和几条理由，而”合适”是一个无法独立检验的判断。使用者只能在没有任何证据的情况下全盘接受或全盘拒绝，协作流程在这里实际上断裂了。

更好的做法是要求 AI 输出可验证的中间产物，而不是直接给出结论：

\begin{quote}
请分别用三种方法——历史均值基准、线性回归和随机森林——在清洗后的数据上训练并预测，输出每种方法在验证集上的 MAE 和 MAPE，并附上实际值与预测值对比的折线图（横轴为时间，纵轴为速度）。我将根据这些结果自行判断哪种方法更适合本项目。
\end{quote}

% ----------------------------------------------------------------
\section{Prompt的结构化管理}

当 Prompt 被用于研究或生产流程时，它应被视为一种需要版本控制和回归测试的轻量程序。提示词、模型版本、采样参数、工具权限和输入数据共同决定输出；只保存提示文本，仍不足以复现实验。提示模式研究将许多有效交互归纳为可复用结构\cite{white2023prompt}，但模式只有在与具体任务的测试样本和验收标准绑定后，才会从经验技巧变成可靠方法。
% ----------------------------------------------------------------

\subsection{把Prompt当代码一样对待}

一次成功的 Prompt 是一种有价值的知识资产。如果只在聊天界面中随意输入并丢弃，它和一段未被保存的代码没有区别——下次需要时无法复用，也无法被团队共享或审查。

把 Prompt 当代码管理，意味着：

\begin{enumerate}
  \item \textbf{保存核心 Prompt：}将用于数据清洗、特征工程、模型训练和结果验证的关键 Prompt 保存为文本文件，与代码放在同一项目目录下，形成文档的一部分；
  \item \textbf{注明意图与约束来源：}在 Prompt 旁边添加简要说明，解释某个约束为何被声明、某段上下文为何被包含——这些说明在日后复查或移交时非常有价值；
  \item \textbf{记录迭代历史：}保存 Prompt 的修改版本，记录每次修改的原因，使分析过程可以被回溯；
  \item \textbf{与团队共享：}在合作项目中，分享关键 Prompt 的方式应与分享代码或分析笔记相同。
\end{enumerate}

\subsection{为常见任务建立模板}

在固定的工作流中，许多任务的结构是相似的，只有具体参数（数据路径、字段名、时间范围等）不同。为这类任务建立 Prompt 模板，可以减少每次重新构建的时间，也有助于保持一致的上下文声明和约束说明。

以下是一个数据探索任务的 Prompt 模板：

\begin{quote}
\textbf{角色：}[描述工作立场和输出受众]

\textbf{数据说明：}数据来源为[描述来源]，文件为[文件名]，时间粒度为[粒度]，关键字段包括：[字段名：含义：单位]。已知数据质量问题：[列出异常类型和处理约定]。

\textbf{任务：}
\begin{enumerate}
  \item 输出数据基本统计摘要（行数、时间范围、每字段非空数量）
  \item 统计每个关键字段的缺失率和异常值比例
  \item 绘制[指定字段]的分布图和时序变化图
  \item 标记任何与预期不符的模式，并说明可能原因
\end{enumerate}

\textbf{约束：}不可修改原始文件；使用 pandas 和 matplotlib；代码包含必要注释。

\textbf{输出格式：}可独立运行的 Python 脚本，图表保存为 PNG 文件，同时输出一个 Markdown 格式的统计摘要表。
\end{quote}

模板的价值不在于机械地填写每个占位符，而在于它提醒你系统地思考每个要素，避免遗漏关键信息。模板本身也应随项目经验迭代更新。

\subsection{跨会话保持上下文}

在持续时间较长的项目中，不同阶段的工作往往分布在多次对话中。AI 不会自动记忆上一次会话的内容，因此在新会话开始时，需要主动重建必要的上下文。

一种有效方式是维护一份“项目上下文摘要”，记录已经确认的关键决策、数据说明和约束，在每次新会话开始时将其附入 Prompt：

\begin{quote}
\textbf{项目上下文摘要}

项目目标：预测某市城区道路 15 分钟后的车速，结果将供信号控制系统参考。

数据：2023 年全年，20 个检测器，5 分钟粒度，字段：timestamp、detector\_id、speed、volume。已完成清洗，输出文件：clean\_detector.csv。

已确认决策：
\begin{itemize}
  \item 速度为零的记录标记为缺失，不保留
  \item 训练集：2023 年 1—10 月；测试集：11—12 月
  \item 所有特征只能使用预测时刻之前的数据
  \item 主评价指标为 MAPE，辅助指标为 MAE
\end{itemize}

当前步骤：构造时间特征和滞后特征。
\end{quote}

摘要不需要包含所有历史决策，只需覆盖当前任务依赖的关键前提。随着项目推进，摘要内容应定期更新，过时的决策应明确标出或删除。

% ----------------------------------------------------------------
\section{实战练习}
% ----------------------------------------------------------------

\subsection*{练习一：诊断并修复问题Prompt}

以下三条 Prompt 各自包含一个或多个问题。请指出问题所在，并将其改写为满足五要素标准的高质量版本。

\begin{enumerate}
  \item“帮我预测明天的交通流量。”
  \item“请用机器学习分析一下哪些因素影响事故发生。”
  \item“清洗数据，然后建模，然后告诉我结果怎么样，模型准不准。”
\end{enumerate}

\subsection*{练习二：从任务背景构建完整Prompt}

从以下任务中选择一个，按照五要素框架（角色、上下文、任务、约束、格式）构建一条完整 Prompt。完成后，检查你的 Prompt 是否覆盖了该任务的所有关键约束。

\begin{enumerate}
  \item 你有一份公交 GPS 轨迹数据，需要识别每辆车每次进站和出站的时间，并统计各站点的平均停靠时长及其分布。
  \item 你有三个月的共享单车借还记录，需要找出借车需求在空间上的热点区域，以便优化车辆调度策略。
  \item 你正在分析一份路段事故数据，需要判断不同路段类型、时段和天气条件下的事故率是否存在显著差异。
\end{enumerate}

\subsection*{练习三：拆解复杂任务}

以下任务目标合理，但不能一次完整地提交给 AI 处理：

\begin{quote}
利用道路检测器数据，建立一个可以预测接下来 30 分钟交通速度的模型；要求能够在新数据到来时更新预测；最终生成一份包含模型说明、性能评估和局限性分析的报告。
\end{quote}

请将它拆解为至少 5 个独立的、可分别验证的子步骤。对每个步骤说明：输入是什么、输出是什么、应如何验证该步骤的结果。

\subsection*{练习四：建立个人Prompt模板库}

结合自己常用的数据分析任务，设计一个可在不同项目中复用的 Prompt 模板。模板应包含角色、上下文占位符、任务结构、需要声明的约束类型，以及输出格式要求。设计完成后，用一个具体任务测试该模板，记录第一次输出中的问题和修正过程。





